% Content imported from Report_newest.docx.md with approved edits.
\chapter{KIẾN TRÚC HỆ THỐNG ĐA THUÊ BAO}\label{v4-chapter-2}

Thiết kế một nền tảng SaaS đa thuê bao đòi hỏi cân bằng nhiều mục tiêu: vừa tiết kiệm chi phí hạ tầng nhờ dùng chung tài nguyên, vừa tránh rò rỉ dữ liệu và hạn chế ảnh hưởng từ hiện tượng hàng xóm ồn ào. Chương~\ref{v4-chapter-2} trình bày hai hướng tổ chức Tách -- Gộp trên mô hình 9 tầng gắn với bốn khía cạnh kiểm soát K1--K4 (Mục~\ref{v4-sec-2-1}), phân tích bốn mô hình Silo, Pool, Bridge, Hybrid cùng cách mở rộng phân vùng theo ô (Mục~\ref{v4-sec-2-2}), và áp dụng khung đánh đổi nhiều yếu tố để chọn mô hình Pool Cấp độ 3 cho ứng dụng Kanban (Mục~\ref{v4-sec-2-3}).

\section{Tổng quan thiết kế hệ thống đa thuê bao}\label{v4-sec-2-1}

Mục này trình bày khung thiết kế hệ thống đa thuê bao dựa trên chiến lược Tách -- Gộp, bốn khía cạnh kiểm soát K1--K4 và Sơ đồ Khung tham chiếu 9 tầng.

\subsection{Chiến lược Tách - Gộp và bốn khía cạnh kiểm soát (K1-K4)}\label{v4-sec-2-1-1}

Thiết kế kiến trúc đa thuê bao xoay quanh ba việc chính: (1) gắn định danh tổ chức vào ngữ cảnh xử lý bằng tenant\_\allowbreak{}id; (2) chọn mức độ chia sẻ hay tách riêng tài nguyên phù hợp ở từng tầng; và (3) thay đổi hành vi phần mềm theo siêu dữ liệu cấu hình trên cùng một bộ mã nguồn.

\textbf{Để biểu diễn cách chia tài nguyên ở từng tầng trong mô hình 9 tầng, nghiên cứu sử dụng hai khái niệm Tách (Dedicated) và Gộp (Shared):}

\begin{itemize}
\item
  \textbf{Tách (Dedicated):} Cấp cho mỗi thuê bao một Thể hiện Thành phần hoặc không gian lưu trữ riêng ở tầng đang xét, chẳng hạn máy chủ riêng, cơ sở dữ liệu riêng hay tiến trình chạy riêng. Cách này giúp cô lập tốt và khoanh vùng lỗi gọn trong một khách hàng, nhưng làm tăng chi phí phần cứng và công sức vận hành.
\item
  \textbf{Gộp (Shared):} Cho nhiều thuê bao dùng chung một Thể hiện Thành phần ở tầng đang xét, ví dụ chung tiến trình ứng dụng, chung bộ nhớ đệm hoặc chung bảng dữ liệu. Cách này giảm chi phí trên mỗi khách hàng nhờ chia sẻ tài nguyên nhàn rỗi, nhưng bắt buộc phải có cơ chế phần mềm để bảo vệ ranh giới giữa các thuê bao.
\end{itemize}

Cách phân tích theo hai hướng Tách và Gộp cho phép phối hợp linh hoạt ranh giới chia sẻ ở từng tầng thay vì phải chọn giống nhau từ trên xuống dưới. Ví dụ, hệ thống có thể Gộp ở Cụm Ứng dụng và Middleware (Tầng 7--9) để dùng chung năng lực tính toán, nhưng Tách lược đồ hoặc cơ sở dữ liệu ở Tầng 5--6 để khách hàng dễ tùy biến cấu trúc bảng. Tương tự, với nền tảng phục vụ nhiều phân khúc khách hàng, nhóm phát triển có thể Gộp dịch vụ phổ thông vào cụm dùng chung nhưng Tách riêng tiến trình xử lý nặng hoặc cụm lưu trữ cho nhóm khách hàng doanh nghiệp.

\textbf{Khi đã chọn Gộp tại bất kỳ tầng nào, hệ thống phải có giải pháp kiểm soát phần mềm tương ứng, chia thành bốn khía cạnh cốt lõi (K1--K4):}

\begin{itemize}
\item
  \textbf{K1 - Kiểm soát định danh và ngữ cảnh thuê bao:} Phân giải chính xác người dùng gắn với tổ chức và duy trì ngữ cảnh hợp lệ xuyên suốt mọi tầng xử lý.
\item
  \textbf{K2 - Cô lập dữ liệu và trạng thái lưu trữ:} Kiểm soát các thao tác đọc, ghi và truy xuất trạng thái trong đúng phạm vi ranh giới của thuê bao.
\item
  \textbf{K3 -- Cô lập hiệu năng và quản trị tài nguyên:} Kiểm soát mức sử dụng CPU, RAM, kết nối cơ sở dữ liệu và tần suất yêu cầu theo từng thuê bao để hạn chế hiện tượng hàng xóm ồn ào.
\item
  \textbf{K4 - Tùy biến cấu hình và hành vi động:} Cho phép điều chỉnh tính năng, quy trình và trường dữ liệu mở rộng theo siêu dữ liệu của từng thuê bao trên cùng một bộ mã nguồn.
\end{itemize}

Nhờ vậy, bài toán kiến trúc được nhìn từ hai phía: nhìn tổng thể xem toàn bộ Thể hiện Hệ thống triển khai theo mô hình nào, và nhìn chi tiết ở từng tầng xem áp dụng Tách hay Gộp cùng cách giải quyết K1--K4 ra sao.

Việc đặt chiến lược Tách -- Gộp trước khi khảo sát các mô hình Silo, Pool, Bridge và Hybrid ở Mục~\ref{v4-sec-2-2} nhằm tránh cách hiểu nhầm rằng đa thuê bao chỉ là chuyện tách hay gộp cơ sở dữ liệu. Đồng thời, đối chiếu lại với ba mục tiêu nền ở Phần I (định danh thuê bao, cô lập dữ liệu và nhất quán trạng thái công việc), phân tích ở Chương~\ref{v4-chapter-1} cho thấy để đạt chuẩn Cấp độ 3 thực sự, hệ thống cần bổ sung thêm cô lập hiệu năng chống hàng xóm ồn ào và tùy biến động qua siêu dữ liệu. Vì vậy, nghiên cứu mở rộng trục kỹ thuật thành bốn khía cạnh K1--K4 xuyên suốt Chương~\ref{v4-chapter-2}--5, còn yêu cầu nhất quán trạng thái công việc được đưa trực tiếp vào giao dịch dữ liệu ở Chương~\ref{v4-chapter-4} bằng khóa bi quan (SELECT ... FOR UPDATE) kết hợp khóa lạc quan (version).

\subsection{Khung kiến trúc tham chiếu và quy ước biểu diễn}\label{v4-sec-2-1-2}

Hình~\ref{fig:v4-2-1} mô tả Khung tham chiếu kiến trúc đa thuê bao 9 tầng, chia rõ Phân vùng Phần mềm Nhận thức Thuê bao (Tầng 4--9) ở phía trên và Phân vùng Hạ tầng Nền tảng (Tầng 1--3) ở phía dưới làm bản đồ áp dụng Tách -- Gộp.

\VfourFigure{2.1}{assets/figures/v4/image9.png}{Khung tham chiếu kiến trúc đa thuê bao 9 tầng.}{fig:v4-2-1}

Các sơ đồ trong báo cáo dùng chung ba quy ước màu và nét vẽ: (1) vùng phần mềm nhận thức thuê bao (Tầng 4--9) vẽ khung nét liền, vùng hạ tầng nền (Tầng 1--3) vẽ khung nét đứt; (2) khung màu riêng đại diện cho không gian hoặc dữ liệu logic của từng thuê bao; (3) khối nét liền thể hiện thành phần Gộp dùng chung, còn khối nét đứt thể hiện thành phần Tách riêng.

\section{Các mô hình triển khai kiến trúc đa thuê bao}\label{v4-sec-2-2}

Trước khi đi vào từng mô hình cụ thể, mục này giải thích nhanh ý nghĩa trực quan của ba tên gọi Silo, Pool và Bridge trong các tài liệu kiến trúc SaaS \cite{v4-18}, \cite{v4-19}:

\begin{itemize}
\item
  Silo (Buồng chứa biệt lập): Mượn từ hình ảnh các tháp chứa ngũ cốc tách rời nhau hoàn toàn để không lẫn nguyên liệu. Trong đa thuê bao, Silo đại diện cho chiến lược Tách: mỗi thuê bao được cấp một vùng tài nguyên riêng biệt (Dedicated).
\item
  Pool (Bể dùng chung): Chỉ việc gom năng lực tính toán, bộ nhớ và lưu trữ vào một bể chung duy nhất. Pool đại diện cho chiến lược Gộp: nhiều thuê bao cùng chia sẻ tài nguyên trong một hệ thống chung.
\item
  Bridge (Cầu nối trung gian): Đóng vai trò trung gian giữa hai thái cực Pool và Silo, tức gộp tài nguyên tính toán ở các tầng trên để tiết kiệm chi phí nhưng tách không gian dữ liệu ở tầng dưới để giữ ranh giới cô lập rõ hơn.
\end{itemize}

Điểm cần nhấn mạnh là Silo, Pool và Bridge áp dụng được cho từng Thể hiện Thành phần ở mỗi tầng, chứ không chỉ riêng ở cơ sở dữ liệu. Chẳng hạn ở Cụm Ứng dụng (Tầng 8--9), hệ thống cũng có thể chọn Silo (mỗi thuê bao có tiến trình riêng, có thể triển khai từ cùng bộ mã nguồn chuẩn), Pool (dùng chung cả mã nguồn và tiến trình thực thi) hoặc Bridge (dùng chung mã nguồn Tầng 9 nhưng tách tiến trình chạy ở Tầng 8). Xét ở mức toàn bộ Thể hiện Hệ thống, có bốn mô hình triển khai chính sau:

\subsection{Mô hình Silo}\label{v4-sec-2-2-1}

Trong mô hình Silo, mỗi Thuê bao sở hữu một Thể hiện Hệ thống dành riêng độc lập từ hạ tầng, ứng dụng đến cơ sở dữ liệu. Toàn bộ các môi trường dành riêng vẫn duy trì một bộ mã nguồn chuẩn thống nhất và được điều phối thông qua Mặt phẳng Điều khiển (Control Plane) tập trung.

\VfourFigure{2.2}{assets/figures/v4/image10.png}{Kiến trúc mô hình Silo.}{fig:v4-2-2}

\subsection{Mô hình Pool}\label{v4-sec-2-2-2}

Mô hình Pool chia sẻ tối đa tài nguyên khi toàn bộ các Thuê bao cùng chạy trên một Thể hiện Hệ thống chung duy nhất từ tầng ứng dụng đến cơ sở dữ liệu. Ranh giới giữa các tổ chức được bảo vệ bằng logic phần mềm gắn với tenant\_\allowbreak{}id xuyên suốt bốn khía cạnh K1--K4.

\VfourFigure{2.3}{assets/figures/v4/image11.png}{Kiến trúc mô hình Pool.}{fig:v4-2-3}

\subsection{Mô hình Bridge}\label{v4-sec-2-2-3}

Bridge là mô hình kết hợp theo chiều dọc: gộp các thành phần tính toán ở Tầng 7--9 (dùng chung ứng dụng và phần trung gian), nhưng tách không gian lưu trữ tại Tầng 5 (cơ sở dữ liệu riêng) hoặc Tầng 6 (lược đồ riêng) cho từng thuê bao \cite{v4-16}.

\VfourFigure{2.4}{assets/figures/v4/image12.png}{Kiến trúc mô hình Bridge.}{fig:v4-2-4}

\Needspace{.55\textheight}
\subsection{Mô hình Hybrid, mở rộng theo ô và so sánh tổng hợp}\label{v4-sec-2-2-4}

Mô hình Hybrid kết hợp đồng thời cả vùng dùng chung (Pool) lẫn vùng dành riêng (Silo) theo phân khúc khách hàng. Nhờ đó, nền tảng giữ chi phí thấp cho nhóm khách hàng phổ thông ở cụm dùng chung, trong khi vẫn đáp ứng yêu cầu bảo mật và SLA cao của doanh nghiệp lớn ở cụm riêng.

\VfourFigure{2.5}{assets/figures/v4/image13.png}{Kiến trúc mô hình Hybrid.}{fig:v4-2-5}

\begin{landscape}
\renewcommand{\VfourImageMaxHeight}{.52\textheight}
Khi số lượng thuê bao vượt quá khả năng phục vụ của một cụm Pool hoặc Hybrid đơn lẻ, hệ thống có thể mở rộng theo chiều ngang bằng kiến trúc Phân vùng theo Ô (Cell-Based Architecture) như trên Hình~\ref{fig:v4-2-6}. Thay vì dồn tất cả khách hàng về một cụm duy nhất, hệ thống đặt một bộ định tuyến toàn cầu ở phía trước để chia tải vào các ô độc lập. Mỗi ô tự tổ chức theo mô hình phù hợp với nhóm khách hàng mình phục vụ (như Ô 01 và Ô 02 chạy Pool cho khách hàng phổ thông, Ô 03 chạy Hybrid đa bậc), giúp giới hạn phạm vi ảnh hưởng khi có sự cố và mở rộng thuận tiện.
\VfourFigure{2.6}{assets/figures/v4/image14.png}{Kiến trúc mở rộng phân vùng theo ô (Cell-Based).}{fig:v4-2-6}

Bảng~\ref{tab:v4-2-1} tóm tắt ưu và nhược điểm của bốn mô hình triển khai trên sáu tiêu chí kiến trúc và kinh tế:
\begingroup\VfourTableSetup
\renewcommand{\thetable}{2.1}
\begin{longtable}{@{}P{\dimexpr 0.18750\linewidth-7.500bp\relax}P{\dimexpr 0.18750\linewidth-7.500bp\relax}P{\dimexpr 0.25000\linewidth-10.000bp\relax}P{\dimexpr 0.18750\linewidth-7.500bp\relax}P{\dimexpr 0.18750\linewidth-7.500bp\relax}@{}}
\caption{So sánh bốn mô hình triển khai kiến trúc đa thuê bao.}\label{tab:v4-2-1}\\
\toprule
Tiêu chí Đối sánh & Mô hình Silo & Mô hình Pool & Mô hình Bridge & Mô hình Hybrid \\
\midrule\endfirsthead
\toprule
Tiêu chí Đối sánh & Mô hình Silo & Mô hình Pool & Mô hình Bridge & Mô hình Hybrid \\
\midrule\endhead
\midrule\endfoot
\bottomrule\endlastfoot
\textbf{1. Tổ chức \emph{thể hiện hệ thống} và \emph{thể hiện thành phần}} & Mỗi \emph{Thuê bao} sở hữu 1 \emph{thể hiện hệ thống dành} riêng (Tách Tầng 1-9; dùng chung bộ mã nguồn và tạo tác triển khai). & Nhiều \emph{Thuê bao} chia sẻ 1 \emph{thể hiện hệ thống dùng chung} duy nhất (Gộp toàn phần Tầng 1-9). & Gộp \emph{các thể hiện thành phần} tầng tính toán (Tầng 7-9); Tách cấu trúc dữ liệu theo cơ sở dữ liệu hoặc lược đồ (Tầng 5-6). & Kết hợp \emph{thể hiện hệ thống dùng chung} cho gói \emph{miễn phí và nâng cao} và \emph{thể hiện hệ thống dành riêng} cho gói \emph{doanh nghiệp}. \\
\textbf{2. Hiệu quả chi phí và kinh tế quy mô} & Chi phí phụ thuộc số môi trường dành riêng và mức tự động hóa. & Chia sẻ công suất nhàn rỗi khi tải và sức chứa cho phép. & Chia sẻ tài nguyên tính toán, duy trì thêm các không gian dữ liệu riêng. & Phối hợp tài nguyên dùng chung và dành riêng theo nhu cầu từng nhóm thuê bao. \\
\textbf{3. Mức độ cô lập và rủi ro hàng xóm ồn ào} & Giảm tranh chấp ở các tài nguyên được cấp riêng. Tài nguyên còn dùng chung vẫn có thể chịu nhiễu. & Cô lập logic bằng K1--K4. Sự cố ở thành phần dùng chung có thể ảnh hưởng nhiều thuê bao. & Dữ liệu tách cấu trúc. Tầng ứng dụng hoặc hệ quản trị cơ sở dữ liệu dùng chung vẫn có thể tranh chấp tài nguyên. & Mức ảnh hưởng tùy theo ranh giới giữa cụm dùng chung và cụm dành riêng. \\
\textbf{4. Độ phức tạp Vận hành và Quản trị} & \textbf{Rất phức tạp khi N lớn:} Phải tự động hóa hạ tầng được mô tả bằng mã để quản lý và nâng cấp N thể hiện hệ thống. & \textbf{Đơn giản nhất về triển khai:} Chỉ vận hành 1 \emph{thể hiện hệ thống}, nhưng đòi hỏi kỷ luật cao trong mã nguồn. & \textbf{Trung bình - Cao:} Triển khai ứng dụng đơn giản, nhưng di trú lược đồ phải lặp qua N lược đồ hoặc cơ sở dữ liệu. & \textbf{Cao nhất về định tuyến:} Đòi hỏi \emph{mặt phẳng điều khiển} và \emph{bộ định tuyến theo gói dịch vụ} điều phối đa hình thái và di trú. \\
\textbf{5. Khả năng mở rộng theo số lượng thuê bao} & Phụ thuộc khả năng cấp phát và quản trị các môi trường dành riêng. & Phụ thuộc sức chứa cụm dùng chung, đặc điểm tải và cơ chế điều tiết. & Phụ thuộc số cơ sở dữ liệu hoặc lược đồ, kết nối và quy trình di trú. & Có thể mở rộng từng nhánh theo đặc điểm tải và yêu cầu cô lập. \\
\textbf{6. Phân khúc Khách hàng Phù hợp Điển hình} & Tài chính, ngân hàng, y tế, chính phủ, tập đoàn yêu cầu tuân thủ pháp lý nghiêm ngặt. & Nền tảng phần mềm dịch vụ (SaaS) phục vụ đa dạng doanh nghiệp và người dùng cá nhân, doanh nghiệp vừa và nhỏ, doanh nghiệp khởi nghiệp và công cụ cộng tác nhóm (quản lý công việc, quản lý quan hệ khách hàng quy mô vừa và nhỏ). & Hệ thống hoạch định nguồn lực doanh nghiệp và quản trị nhân sự tầm trung cần tùy biến bảng dữ liệu hoặc trích xuất/sao lưu cơ sở dữ liệu riêng. & Các nền tảng phần mềm dưới dạng dịch vụ thương mại trưởng thành có bảng giá đa bậc. \\
\end{longtable}
\endgroup
\end{landscape}

\section{Khung đánh giá đa yếu tố và lựa chọn mô hình cho đề tài}\label{v4-sec-2-3}

Bốn mô hình Silo, Bridge, Hybrid và Pool đều có điểm mạnh riêng và không có một mô hình nào đúng cho mọi hoàn cảnh. Khi lựa chọn kiến trúc, người thiết kế cần cân nhắc đồng thời bốn yếu tố sau:

\begin{enumerate}
\def\labelenumi{\arabic{enumi}.}
\item
  \textbf{Thị trường mục tiêu, quy mô khách hàng và chi phí:} Silo phù hợp khi khách hàng cần tài nguyên riêng và có thể chịu chi phí tương ứng. Pool có thể giảm chi phí biên trong phạm vi tài nguyên dùng chung còn đủ. Hybrid có thể phục vụ các nhóm khách hàng có nhu cầu và ngân sách khác nhau.
\item
  \textbf{Yêu cầu pháp lý, mức độ cô lập và khả năng chịu nhiễu:} Cần phân biệt rõ trường hợp bắt buộc cách ly phần cứng hoặc khóa mã hóa riêng theo luật định (buộc dùng Silo hoặc nhánh riêng của Hybrid) với trường hợp cho phép cô lập bằng logic có kiểm chứng chặt chẽ trong phần mềm và cơ sở dữ liệu (Pool). Nếu chỉ cần tách không gian sao lưu và cấu trúc bảng theo từng khách hàng thì Bridge là lựa chọn trung gian phù hợp.
\item
  \textbf{Mức độ tùy biến và cách nâng cấp lược đồ dữ liệu:} Nếu khách hàng cần tự thêm bảng vật lý hoặc chỉnh sửa trực tiếp cấu trúc bảng bằng câu lệnh DDL, Bridge hoặc Silo đáp ứng tự nhiên hơn. Ngược lại, nếu sự khác biệt về giao diện, quy trình và trường dữ liệu có thể mô tả bằng siêu dữ liệu trên cùng một lược đồ bảng chung, Pool sẽ giữ tốc độ nâng cấp ở mức O(1).
\item
  \textbf{Sự đánh đổi công sức giữa vận hành hạ tầng và lập trình phần mềm:} Dùng chung một bộ mã nguồn ở Tầng 9 không có nghĩa là dùng chung tiến trình chạy. Silo ít rủi ro lẫn dữ liệu trong mã nguồn nhưng đẩy gánh nặng sang tự động hóa và quản trị N cụm máy chủ; Bridge phát sinh việc chạy di trú lặp lại qua N lược đồ hay cơ sở dữ liệu; Hybrid tăng độ phức tạp định tuyến ở cổng biên; còn Pool gọn nhẹ nhất khi triển khai nhưng dồn trách nhiệm an toàn và kiểm soát K1--K4 vào kỷ luật thiết kế phần mềm.
\end{enumerate}

\begin{landscape}
\renewcommand{\VfourImageMaxHeight}{.72\textheight}
Vì thế, trước khi chọn mô hình tổ chức tài nguyên, nhóm phát triển cần nắm rõ đối tượng khách hàng, kỳ vọng chi phí, ràng buộc bảo mật và mức độ tùy biến nghiệp vụ thực tế. Khi đặt cả bốn yếu tố này lên cùng một khung đánh giá, quyết định Tách -- Gộp ở từng tầng và hướng mở rộng về sau sẽ trở nên rõ ràng hơn.
\VfourFigure{2.7}{assets/figures/v4/image15.png}{Khung đánh đổi đa yếu tố khi lựa chọn mô hình đa thuê bao.}{fig:v4-2-7}
\end{landscape}

\textbf{Áp dụng bốn yếu tố đánh giá trên vào bài toán quản lý công việc Kanban của đề tài (Hình~\ref{fig:v4-2-8}), từng lựa chọn được đối chiếu như sau:}

\begin{enumerate}
\def\labelenumi{\arabic{enumi}.}
\item
  Về kinh tế và quy mô: Nhóm người dùng mục tiêu của ứng dụng là số đông doanh nghiệp vừa và nhỏ, phòng ban và nhóm dự án cần đăng ký dùng ngay với chi phí thấp và thời gian làm việc tập trung theo giờ hành chính, nên đề tài ưu tiên chia sẻ hạ tầng theo giả định chi phí ở Mục~\ref{v4-sec-1-1-2}.
\item
  Về pháp lý và cô lập: Phạm vi nguyên mẫu giả định có thể dùng cơ chế cô lập logic ở tầng ứng dụng và cơ sở dữ liệu. Yêu cầu tách hạ tầng vật lý hoặc tuân thủ trong một triển khai cụ thể cần được đánh giá theo tổ chức tiếp nhận.
\item
  Về tùy biến và lược đồ: Mọi thuê bao đều dùng chung khung dữ liệu cốt lõi (Bảng -- Cột -- Thẻ), còn phần khác biệt về giao diện, giới hạn WIP và trường động đều lưu được qua siêu dữ liệu trên một lược đồ chung, không cần gánh độ phức tạp bảo trì N lược đồ của Bridge.
\item
  Về cân bằng vận hành và kỹ thuật: Bắt đầu bằng một Thể hiện Hệ thống tập trung (chưa cần dựng bộ định tuyến đa vùng phức tạp của Hybrid ở bước đầu) giúp giảm công sức triển khai và tập trung nguồn lực giải quyết trọn vẹn K1--K4.
\end{enumerate}

\VfourFigure{2.8}{assets/figures/v4/image16.png}{Áp dụng khung đánh giá để lựa chọn mô hình Pool cho ứng dụng Kanban.}{fig:v4-2-8}

Từ các phân tích trên, đề tài chọn Mô hình Pool ở Cấp độ trưởng thành 3 làm kiến trúc mục tiêu cho ứng dụng SaaS Kanban. Toàn bộ các thành phần từ Tầng 1 đến Tầng 9 đều tổ chức theo hướng gộp dùng chung. Tuy nhiên, lựa chọn Pool xác định cách chia sẻ tài nguyên trong phạm vi giả định; khi gỡ bỏ vách ngăn phần cứng, hệ thống phải đối mặt trực tiếp với nguy cơ lẫn ngữ cảnh, rò rỉ dữ liệu và tranh chấp tải. Chương~\ref{v4-chapter-3} sẽ đi vào từng cụm tầng để xử lý cụ thể bốn khía cạnh kiểm soát K1--K4.

Việc chọn Pool Cấp độ 3 dựa trên nhận thức rõ ràng về sự đánh đổi: Silo và Bridge có sẵn vách ngăn phần cứng hoặc lược đồ tự nhiên mạnh hơn, nhưng kéo theo công việc cấp phát hoặc di trú trên nhiều tài nguyên riêng. Trong khi đó, Pool Cấp độ 3 vừa hợp với bài toán kinh tế của ứng dụng Kanban vừa là môi trường tốt nhất để kiểm nghiệm trọn vẹn các kỹ thuật phần mềm K1--K4. Ngoài ra, so với Mục 6 Phần I, việc so sánh Silo, Bridge, Pool và Hybrid được thực hiện ở mức phân tích kiến trúc tại Chương~\ref{v4-chapter-2}--3, còn phần cài đặt Ứng dụng và đo kiểm định lượng ở Chương~\ref{v4-chapter-4}--5 đi sâu vào Pool Cấp độ 3 qua đối chứng hai cấu hình cơ sở (Baseline) và đề xuất (Proposed).
